\chapter{Arithmétique : divisibilité, congruences et numération}

\noindent\textit{L'arithmétique étudie les nombres entiers et leurs propriétés. Discipline
reine de la série C, elle développe la rigueur du raisonnement (récurrence, disjonction
de cas, raisonnement par l'absurde) et trouve des applications modernes en cryptographie
et en informatique. Ce premier chapitre pose les fondations : divisibilité, division
euclidienne, congruences et écriture des nombres dans différentes bases.}
\vspace{8pt}

% ═══════════════════════════════════════════════════════════════
\section{Divisibilité dans $\Z$}

\subsection{Diviseurs et multiples}

\begin{definition}[Divisibilité]
Soient $a$ et $b$ deux entiers relatifs. On dit que $b$ \textbf{divise} $a$ (ou que $a$ est un \textbf{multiple} de $b$), et l'on note $b\mid a$, s'il existe un entier relatif $k$ tel que $a=k\,b$.
\end{definition}

\begin{exemple}
$7\mid 42$ car $42=6\times 7$ ; $\;-3\mid 12$ car $12=(-4)\times(-3)$. Tout entier divise $0$ (car $0=0\times b$), et $1$ divise tout entier.
\end{exemple}

\begin{propriete}[Règles de calcul]
Pour tous entiers $a,b,c$ :
\begin{itemize}
\item si $a\mid b$ et $b\mid c$, alors $a\mid c$ \quad(\emph{transitivité}) ;
\item si $a\mid b$ et $a\mid c$, alors $a\mid (bu+cv)$ pour tous entiers $u,v$ \quad(\emph{combinaison linéaire}) ;
\item si $a\mid b$ et $b\mid a$, alors $a=b$ ou $a=-b$ ;
\item si $a\mid b$ et $b\neq 0$, alors $|a|\le |b|$.
\end{itemize}
\end{propriete}

\begin{remarque}
La propriété de combinaison linéaire est un outil \emph{fondamental} : pour montrer qu'un entier $d$ divise une expression, on cherche à l'écrire comme combinaison d'entiers que $d$ divise déjà.
\end{remarque}

\subsection{La division euclidienne}

\begin{theoreme}[Division euclidienne dans $\N$ puis $\Z$]
Pour tout entier $a$ et tout entier $b>0$, il existe un \textbf{unique} couple d'entiers $(q,r)$ tel que
\[
a=b\,q+r \qquad\text{avec}\qquad 0\le r<b.
\]
$q$ est le \textbf{quotient} et $r$ le \textbf{reste} de la division euclidienne de $a$ par $b$.
\end{theoreme}

\begin{illus}
\begin{tikzpicture}[x=1cm,>=Stealth]
\draw[->] (-0.4,0) -- (10.6,0);
\foreach \i/\lab in {0/0,2/b,4/2b,6/3b,8/4b}{
  \draw (\i,0.12) -- (\i,-0.12) node[below=2pt]{\small $\lab$};}
\draw[onbo,line width=1.2pt] (8,0.12) -- (8,-0.12);
\node[onbo] at (8,0.5) {\small $qb$};
\filldraw[onbo] (9.1,0) circle (2.2pt) node[above=4pt]{\small $a$};
\draw[onbgreen,decorate,decoration={brace,amplitude=5pt},line width=1pt] (9.1,-0.5) -- (8,-0.5)
  node[midway,below=6pt,onbgreen]{\small $r$};
\node at (4,1.1) {\small $a=bq+r,\quad 0\le r<b$ : $a$ se situe entre $qb$ et $(q{+}1)b$.};
\end{tikzpicture}
\fig{Figure 1.1 — Interprétation de la division euclidienne sur l'axe des entiers.}
\end{illus}

\begin{exemple}
Division de $47$ par $6$ : $47=6\times 7+5$, donc $q=7$ et $r=5$.
Pour un entier \emph{négatif}, on garde $0\le r<b$ : $-47=6\times(-8)+1$, donc le reste est $1$ (et non $-5$).
\end{exemple}

\subsection{Critères de divisibilité}

\begin{propriete}[Critères usuels]
Un entier écrit en base dix est divisible par :
\begin{itemize}
\item \textbf{2} (resp. \textbf{5}) si son chiffre des unités est divisible par $2$ (resp. par $5$) ;
\item \textbf{4} si le nombre formé par ses deux derniers chiffres l'est ;
\item \textbf{3} (resp. \textbf{9}) si la somme de ses chiffres est divisible par $3$ (resp. par $9$) ;
\item \textbf{11} si la somme alternée de ses chiffres est divisible par $11$.
\end{itemize}
\end{propriete}

% ═══════════════════════════════════════════════════════════════
\section{Congruences modulo $n$}

\begin{definition}[Congruence]
Soit $n$ un entier naturel non nul. Deux entiers $a$ et $b$ sont \textbf{congrus modulo} $n$, ce que l'on note $a\equiv b\;[n]$, lorsque $n\mid (a-b)$, c'est-à-dire lorsque $a$ et $b$ ont le \textbf{même reste} dans la division euclidienne par $n$.
\end{definition}

\begin{illus}
\begin{tikzpicture}[>=Stealth,scale=1]
\def\R{1.5}
\foreach \i/\ang in {0/90,1/18,2/-54,3/-126,4/162}{
  \node[circle,draw=onbo,fill=onbsoft,inner sep=2.5pt] (n\i) at (\ang:\R) {\small $\i$};}
\foreach \i/\j in {0/1,1/2,2/3,3/4,4/0}{
  \draw[->,onbo,line width=0.9pt] (n\i) to[bend left=18] (n\j);}
\node at (0,0) {\small $\bmod\,5$};
\node[right=1.6cm,align=left,text width=5.4cm] at (\R,0)
{\small La congruence modulo $5$ « enroule » $\Z$ sur $5$ classes. Ajouter $1$ fait avancer d'un cran : $\dots\equiv 7\equiv 12\equiv 2\;[5]$.};
\end{tikzpicture}
\fig{Figure 1.2 — Roue des congruences modulo $5$ (les cinq classes de restes).}
\end{illus}

\begin{theoreme}[Compatibilité avec les opérations]
Si $a\equiv b\;[n]$ et $c\equiv d\;[n]$, alors :
\[
a+c\equiv b+d\;[n],\qquad a\,c\equiv b\,d\;[n],\qquad a^{p}\equiv b^{p}\;[n]\ \ (p\in\N).
\]
\end{theoreme}

\begin{aretenir}
Les congruences sont l'outil idéal pour déterminer un \textbf{reste} ou le \textbf{chiffre des unités} d'un grand nombre, et pour étudier la \textbf{divisibilité} d'une expression dépendant de $n$ : on remplace chaque entier par un représentant simple de sa classe.
\end{aretenir}

\begin{exemple}
Reste de $2^{2024}$ modulo $7$. On a $2^{3}=8\equiv 1\;[7]$. Comme $2024=3\times 674+2$,
$2^{2024}=(2^{3})^{674}\times 2^{2}\equiv 1^{674}\times 4\equiv 4\;[7]$. Le reste est $\mathbf{4}$.
\end{exemple}

% ═══════════════════════════════════════════════════════════════
\section{Nombres premiers et décomposition}

\begin{definition}[Nombre premier]
Un entier $p\ge 2$ est \textbf{premier} s'il n'admet que deux diviseurs positifs : $1$ et lui-même. Un entier $\ge 2$ non premier est dit \textbf{composé}.
\end{definition}

\begin{theoreme}[Décomposition en facteurs premiers]
Tout entier $n\ge 2$ s'écrit de manière \textbf{unique} (à l'ordre des facteurs près) comme produit de nombres premiers :
\[
n=p_{1}^{\alpha_{1}}\,p_{2}^{\alpha_{2}}\cdots p_{k}^{\alpha_{k}}.
\]
\end{theoreme}

\begin{illus}
\begin{tikzpicture}[>=Stealth,level distance=10mm,sibling distance=20mm,
  every node/.style={font=\small},
  pf/.style={circle,draw=onbo,fill=onbsoft,inner sep=1.5pt},
  cp/.style={rectangle,draw=onbink,fill=white,inner sep=2pt}]
\node[cp]{$360$}
  child{node[pf]{$2$}}
  child{node[cp]{$180$}
    child{node[pf]{$2$}}
    child{node[cp]{$90$}
      child{node[pf]{$2$}}
      child{node[cp]{$45$}
        child{node[pf]{$3$}}
        child{node[cp]{$15$}
          child{node[pf]{$3$}}
          child{node[pf]{$5$}}}}}};
\node[right=3.2cm,align=left,text width=4.6cm] at (0,-2)
{\small On divise successivement par les nombres premiers croissants :
\[360=2^{3}\times 3^{2}\times 5.\]};
\end{tikzpicture}
\fig{Figure 1.3 — Arbre de décomposition de $360$ en facteurs premiers.}
\end{illus}

\begin{propriete}[Nombre de diviseurs]
Si $n=p_{1}^{\alpha_{1}}\cdots p_{k}^{\alpha_{k}}$, alors le nombre de diviseurs positifs de $n$ est
\[
(\alpha_{1}+1)(\alpha_{2}+1)\cdots(\alpha_{k}+1).
\]
\end{propriete}

\begin{exemple}
$360=2^{3}\times 3^{2}\times 5^{1}$ possède $(3{+}1)(2{+}1)(1{+}1)=24$ diviseurs positifs.
\end{exemple}

% ═══════════════════════════════════════════════════════════════
\section{Systèmes de numération}

\begin{definition}[Écriture en base $b$]
Soit $b\ge 2$ un entier. Tout entier $N\ge 1$ s'écrit de façon unique
\[
N=a_{p}\,b^{p}+a_{p-1}\,b^{p-1}+\dots+a_{1}\,b+a_{0},\qquad 0\le a_{i}<b,\ a_{p}\neq 0,
\]
que l'on note $N=\overline{a_{p}a_{p-1}\dots a_{0}}^{\,b}$. Les $a_i$ sont les \textbf{chiffres} de $N$ en base $b$.
\end{definition}

\begin{exemple}
$\overline{1011}^{\,2}=1\cdot 2^{3}+0\cdot 2^{2}+1\cdot 2+1=11$. Réciproquement, pour écrire $45$ en base $2$, on effectue des divisions successives par $2$ et l'on lit les restes de bas en haut : $45=\overline{101101}^{\,2}$.
\end{exemple}

% ═══════════════════════════════════════════════════════════════
\section{Méthodes \& savoir-faire}

\methode{Déterminer un reste à l'aide des congruences}
Remplacer chaque facteur par un représentant simple de sa classe modulo $n$, puis exploiter les puissances qui valent $1$.
\begin{exemple}
Chiffre des unités de $7^{2025}$ : c'est le reste modulo $10$. Or $7^{1}\equiv 7$, $7^{2}\equiv 9$, $7^{3}\equiv 3$, $7^{4}\equiv 1\;[10]$ : le cycle est de longueur $4$. Comme $2025=4\times 506+1$, $7^{2025}\equiv 7^{1}\equiv 7\;[10]$. Le chiffre des unités est $\mathbf 7$.
\end{exemple}

\methode{Démontrer une divisibilité par récurrence}
Pour montrer que $n\mid u_k$ pour tout $k$, on initialise puis on suppose $n\mid u_k$ et on écrit $u_{k+1}$ comme combinaison faisant apparaître $u_k$.
\begin{exemple}
Montrons que $3\mid 4^{k}-1$ pour tout $k\in\N$. \emph{Initialisation :} $4^{0}-1=0$, divisible par $3$. \emph{Hérédité :} si $3\mid 4^{k}-1$, alors $4^{k+1}-1=4\cdot 4^{k}-1=4(4^{k}-1)+3$ ; somme de deux multiples de $3$, donc $3\mid 4^{k+1}-1$.
\end{exemple}

\methode{Décomposer en facteurs premiers et compter les diviseurs}
Diviser successivement par les premiers croissants, puis appliquer la formule du produit des $(\alpha_i+1)$.

\methode{Changer de base}
\emph{Vers la base dix} : développer les puissances. \emph{Depuis la base dix} : divisions euclidiennes successives par $b$, les restes lus de bas en haut donnent les chiffres.

% ═══════════════════════════════════════════════════════════════
\section{Exercices}

\rubrique{Divisibilité et division euclidienne}
\exo{}\diff{1} Effectuer la division euclidienne de $2024$ par $7$, puis de $-2024$ par $7$.

\exo{}\diff{1} Déterminer tous les diviseurs positifs de $36$.

\exo{}\diff{2} Montrer que pour tout entier $n$, le produit $n(n+1)$ est divisible par $2$, et $n(n+1)(n+2)$ par $6$.

\rubrique{Congruences}
\exo{}\diff{1} Déterminer le reste de la division de $5^{100}$ par $4$.

\exo{}\diff{2} Quel est le chiffre des unités de $3^{2025}$ ?

\exo{}\diff{3} Montrer que pour tout entier $n$, $7\mid 3^{2n+1}+2^{n+2}$.

\rubrique{Nombres premiers et numération}
\exo{}\diff{2} Décomposer $1\,800$ en facteurs premiers et donner son nombre de diviseurs.

\exo{}\diff{2} Écrire $45$ en base $2$ et $171$ en base $16$. Convertir $\overline{2301}^{\,4}$ en base dix.

\rubrique{Problème type Baccalauréat}
\exo{}\diff{3} On pose $A=n^{3}-n$ pour $n\in\N$.
\begin{enumerate}[label=\arabic*)]
\item Factoriser $A$.
\item Montrer que $A$ est divisible par $6$ pour tout entier $n$.
\item En déduire le reste de la division de $n^{3}$ par $6$ en fonction de celui de $n$.
\end{enumerate}

% ═══════════════════════════════════════════════════════════════
\section{Corrigés}

\corrige{de l'exercice 1}
$2024=7\times 289+1$ (car $7\times 289=2023$), donc $q=289$, $r=1$. Pour $-2024$ : on veut $0\le r<7$ ; $-2024=7\times(-290)+6$ (car $7\times(-290)=-2030$ et $-2030+6=-2024$), donc $q=-290$, $r=6$.

\corrige{de l'exercice 2}
$36=2^{2}\times 3^{2}$. Ses diviseurs positifs sont $1,2,3,4,6,9,12,18,36$ (soit $(2{+}1)(2{+}1)=9$ diviseurs).

\corrige{de l'exercice 3}
Parmi deux entiers consécutifs $n,n+1$, l'un est pair : $2\mid n(n+1)$. Parmi trois entiers consécutifs $n,n+1,n+2$, l'un est multiple de $3$ et l'un est pair : le produit est divisible par $2$ et par $3$, donc par $6$.

\corrige{de l'exercice 4}
$5\equiv 1\;[4]$, donc $5^{100}\equiv 1^{100}\equiv 1\;[4]$ : le reste est $\mathbf 1$.

\corrige{de l'exercice 5}
Les puissances de $3$ modulo $10$ suivent le cycle $3,9,7,1$ de longueur $4$. Comme $2025=4\times 506+1$, $3^{2025}\equiv 3^{1}\equiv 3\;[10]$ : le chiffre des unités est $\mathbf 3$.

\corrige{de l'exercice 6}
Modulo $7$ : $3^{2}=9\equiv 2$, donc $3^{2n+1}=3\cdot(3^{2})^{n}\equiv 3\cdot 2^{n}\;[7]$. Ainsi $3^{2n+1}+2^{n+2}\equiv 3\cdot 2^{n}+4\cdot 2^{n}=7\cdot 2^{n}\equiv 0\;[7]$. Donc $7\mid 3^{2n+1}+2^{n+2}$.

\corrige{de l'exercice 7}
$1800=2^{3}\times 3^{2}\times 5^{2}$ ; nombre de diviseurs : $(3{+}1)(2{+}1)(2{+}1)=36$.

\corrige{de l'exercice 8}
$45=\overline{101101}^{\,2}$ (divisions par $2$). $171=16\times 10+11=\overline{\text{AB}}^{\,16}$ avec $10=\text A$, $11=\text B$, soit $\overline{\text{AB}}^{\,16}$. Enfin $\overline{2301}^{\,4}=2\cdot 4^{3}+3\cdot 4^{2}+0\cdot 4+1=128+48+0+1=177$.

\corrige{du problème}
1) $A=n^{3}-n=n(n^{2}-1)=n(n-1)(n+1)=(n-1)\,n\,(n+1)$ : produit de trois entiers consécutifs.\;
2) Parmi trois entiers consécutifs, il y a au moins un multiple de $2$ et un multiple de $3$ ; $2$ et $3$ étant premiers entre eux, $A$ est divisible par $6$.\;
3) De $n^{3}-n\equiv 0\;[6]$ on tire $n^{3}\equiv n\;[6]$ : le reste de $n^{3}$ modulo $6$ est \textbf{égal} à celui de $n$.
